\begin{textsample}{NEG Dim 1 – 2021\_c4 – Score -5.00 – 2021\_c4/enade\_2021\_art\_vis\_426539035c.txt}  \label{ex:f1_neg_069}
CAPÍTULO 4 PERCEPÇÃO DA PROVA Nas análises feitas neste capítulo, constam as percepções dos concluintes da Área de Artes Visuais ( Licenciatura ) sobre a prova aplicada no Enade/2021.
Essas percepções foram mensuradas por meio de nove questões que avaliaram desde o grau de dificuldade da prova até o tempo gasto para concluí -la.
As percepções sobre a prova foram relacionadas ao desempenho dos estudantes, à Grande Região de funcionamento do curso, à Categoria Administrativa e ao tipo de Organização Acadêmica da IES.
O questionário de percepção da prova encontra -se ao final do Anexo VIII, onde \textbf{está} a reprodução do exame.
O desempenho dos estudantes foi classificado em quatro quartos.
Para tanto, esse desempenho foi ordenado de forma ascendente.
O percentil 25, P25, também conhecido como primeiro quartil, é a nota de desempenho que deixa um quarto ( 25 \% ) dos valores observados abaixo e três quartos acima.
Na Figura 1, apresenta -se uma ilustração desse conceito.
O quarto inferior de desempenho é composto pelas notas abaixo do primeiro quartil.
Já o percentil 75, P75, também conhecido como terceiro quartil, é o valor para o qual há três quartos ( 75 \% ) dos dados abaixo e um quarto acima dele.
O quarto superior de desempenho é composto pelas notas iguais ou acima do terceiro quartil.
O percentil 50, P50, também conhecido como mediana, é o valor que divide as notas em dois conjuntos de igual tamanho.
O segundo quarto inclui valores entre o primeiro quartil ( P25 ) e a mediana.
O terceiro quarto contém os valores entre a mediana ( P50 ) e o terceiro quartil ( P75 ).
Vale ressaltar que percentis, quartis e medianas são pontos que, não obrigatoriamente, pertencem ao conjunto original de dados, ao passo que os quartos são subconjuntos dos dados originais.
A seguir, serão apresentados gráficos com resultados selecionados, relativos às nove questões avaliadas por grupos de estudantes.
Nas barras dos gráficos, apresenta -se o percentual de alunos que assinalou uma das opções ou a soma das porcentagens daqueles que assinalaram duas ( ou três ) delas.
Por exemplo, para as questões 1 e 2, nos gráficos, é apresentada a porcentagem total de participantes que assinalou as opções ( D ) Difícil e ( E ) Muito difícil.
Em \textbf{cada} barra, foram assinalados, também, os extremos do intervalo de confiança de 95 \% como linhas verticais unidas por uma linha horizontal na forma da letra H maiúscula.
O estimador de um parâmetro com um \textbf{certo} nível de confiança ( e.g. 95,0 \% ) deve conter o parâmetro no intervalo de confiança em 95 \% das vezes.
Na comparação entre os estimadores dos parâmetros de duas classes de uma dada categoria ( e.g.
Norte e Nordeste nas Grandes Regiões, de primeiro e último quarto dentro de desempenho, de IES Públicas e Privada, ou de Universidades e Faculdades ), associados aos seus respectivos intervalos de confiança, diz -se que não há diferença estatisticamente significativa entre os parâmetros das duas categorias quando há interseção entre os intervalos de confiança e que há diferença, se os intervalos de confiança são disjuntos ( para maiores informações vide Glossário ).
Nas Tabelas do Anexo II, são apresentados os valores absolutos e a distribuição percentual18 das alternativas válidas das nove questões, segundo o mesmo recorte de desempenho dos alunos e da Grande Região de funcionamento do curso, de Categoria Administrativa e do tipo de Organização Acadêmica da IES. 4.1 GRAU DE DIFICULDADE DA PROVA 4.1.1 COMPONENTE DE FORMAÇÃO GERAL Ao avaliarem “ Qual o grau de dificuldade desta prova na parte de Formação Geral? ” ( Questão 1 ), 36,2 \% do grupo de inscritos e presentes optaram pelas alternativas Difícil ou Muito difícil.
Entretanto, para mais da metade dos estudantes ( 54,7 \% ), o Componente de Formação Geral da prova foi avaliado com grau de dificuldade Médio ( Gráfico 4.1, Gráfico 4.2 e, no Anexo II, a Tabela II.1 ).
O percentual de estudantes que considerou a prova Difícil ou Muito difícil foi maior na região Norte, onde a proporção foi de 45,0 \%, enquanto a de menor proporção foi a Nordeste, com 23,1 \%.
No Gráfico 4.1, é possível observar que as diferenças entre as regiões Nordeste e as regiões Norte, Sudeste e Sul são estatisticamente significativas, assim como a diferença entre a região Centro-Oeste e as regiões Norte, Sudeste e Sul.
Nas Grandes Regiões, a proporção de presentes à prova que avaliou o Componente de Formação Geral como de grau de dificuldade Médio \textbf{esteve} entre 47,9 \%, na região Norte, e 65,8 \%, na região Centro-Oeste. 18 Cumpre \textbf{lembrar} uma das convenções para tabelas numéricas ( ao final do Sumário ) sobre a possibilidade de a soma das partes não resultar em 100 \% por questões de arredondamento.
O percentual de alunos que considerou o componente de Formação Geral da prova Difícil ou Muito difícil foi decrescente em função dos quartos de desempenho : 52,6 \%, no primeiro quarto, e 17,3 \%, no quarto, o grupo de melhor desempenho na prova.
Nos quartos de desempenho intermediários, a proporção de alunos que considerou a prova Difícil ou Muito difícil foi de 44,3 \%, no segundo quarto, e de 31,7 \%, no terceiro quarto.
As diferenças entre os quartos de desempenho são estatisticamente significativas.
Para todos os quartos de desempenho, a alternativa modal para essa pergunta foi a Médio, com 40,8 \% e 66,2 \% dos respondentes nos quartos extremos, primeiro e quarto, respectivamente.
Estas proporções são crescentes com o desempenho.
O percentual de estudantes que considerou o componente de Formação Geral da prova Difícil ou Muito difícil foi maior entre os estudantes de Centros Universitários, alcançando 47,3 \%, enquanto entre os estudantes das Faculdades, a Organização Acadêmica com a menor incidência, alcançou 16,7 \%.
No Gráfico 4.3, é possível observar que as diferenças entre Universidades, Centros Universitários e Faculdades são estatisticamente significativas.
Nas Organizações Acadêmicas, a proporção de presentes à prova que avaliou o Componente de Formação Geral como de grau de dificuldade Médio \textbf{esteve} entre 46,7 \%, nos CEFET/IF, e 70,0 \% nas Faculdades ( ver também a Tabela II.2, no Anexo II ).
O percentual de estudantes que considerou o componente de Formação Geral da prova Difícil ou Muito difícil foi maior entre os estudantes de cursos de IES Privadas ( 44,0 \% ), com diferença estatisticamente significativa da proporção de IES Públicas ( 23,5 \% ).
Já a proporção de presentes à prova que avaliou o Componente de Formação Geral como de grau de dificuldade Médio foi de 63,3 \% para os de IES Públicas e de 49,5 \% para os de IES Privadas ( ver Gráfico 4.4 ). 4.1.2 COMPONENTE DE CONHECIMENTO ESPECÍFICO Ao responderem à Questão 2 – “ Qual o grau de dificuldade desta prova na parte de Componente Específico? ” – 38,8 \% do grupo de estudantes classificaram -na como Difícil ou Muito difícil.
Além disso, o Componente modal de Conhecimento Específico da prova foi considerado de grau de dificuldade Médio por 53,8 \% dos alunos ( Gráfico 4.5, Gráfico 4.6, e, no Anexo II, a Tabela II.3 ).
A análise das respostas dos estudantes quanto ao grau de dificuldade do Componente de Conhecimento Específico da prova, agregado por Grande Região, mostra diferenças estatisticamente significativas entre proporções de alunos que a avaliaram como Difícil ou Muito difícil.
Tais diferenças são observadas entre a região com a maior proporção, a Norte ( 42,9 \% ), e a região com a menor proporção, a Nordeste ( 28,3 \% ).
As diferenças entre a região Nordeste e as regiões Sudeste ( 40,7 ) e Sul ( 39,2 \% ) também são estatisticamente significativas.
O percentual de alunos que classificou o grau de dificuldade como Médio, no Componente de Conhecimento Específico, variou de 51,4 \% a 61,4 \%, nas regiões Norte e Nordeste, respectivamente.
Considerando -se a avaliação da dificuldade das questões do Componente de Conhecimento Específico da prova de acordo com o desempenho dos estudantes, observa -se uma tendência decrescente dos que consideraram a parte específica Difícil ou Muito difícil, variando de 25,1 \% ( último quarto ) a 48,6 \% ( primeiro quarto ).
As demais proporções foram de 46,7 \% e 35,5 \%, respectivamente, no segundo e terceiro quartos.
Observa -se que apenas não há diferença estatisticamente significativa dos resultados entre o primeiro e o segundo quartos.
Já os que responderam que o grau de dificuldade das questões do Componente de Conhecimento Específico da prova foi Médio atingiram percentagens que variaram de 46,0 \%, no primeiro quarto, a 62,5 \%, no último, valores crescentes com o desempenho.
Considerando -se a avaliação da dificuldade das questões do Componente de Conhecimento Específico da prova, o percentual de estudantes que as considerou Difícil ou Muito difícil foi maior entre os estudantes de Centros Universitários, atingindo 47,0 \% ; já entre os de Faculdade, a Organização Acadêmica com a menor incidência, essa proporção foi de 26,7 \%.
No Gráfico 4.7, é possível observar que a diferença entre Universidades ( 34,6 \% ) e Centros Universitários é estatisticamente significativa, bem como a entre Centros Universitários e Faculdades.
Nas Organizações Acadêmicas, a proporção de presentes à prova que considerou o Componente de Conhecimento Específico de grau de dificuldade Médio \textbf{esteve} entre 48,0 \%, nos Centros Universitários, e 63,3 \%, nas Faculdades ( ver também Tabela II.4, no Anexo II ).
O percentual de estudantes que considerou o Componente Específico da prova Difícil ou Muito difícil foi maior entre os estudantes de cursos de IES Privadas ( 46,0 \% ), com diferença estatisticamente significativa da proporção em IES Públicas ( 27,1 \% ).
Já a proporção de presentes à prova que avaliou este componente como de grau de dificuldade Médio foi de 62,0 \% para os de IES Públicas e de 48,7 \% para os de IES Privadas ( ver Gráfico 4.8 e Tabela II.4, no Anexo II ). 4.2 EXTENSÃO DA PROVA EM RELAÇÃO AO TEMPO TOTAL Indagados quanto à extensão da prova, em relação ao tempo total oferecido para a sua resolução ( Questão 3 ), os estudantes apontaram, com maior incidência, a alternativa que considerava a extensão Adequada, para todas as agregações consideradas ( Gráfico 4.9, Gráfico 4.10, e, no Anexo II, a Tabela II.5 ).
O percentual de alunos que respondeu ser a extensão da prova Adequada foi de 45,0 \%.
Já 45,6 \% dos inscritos presentes consideraram que a prova foi Longa ou Muito longa, e 9,4 \% a avaliaram -na como Curta ou Muito curta.
Entre as Grandes Regiões, os que avaliaram a prova como Longa ou Muito longa, em relação ao tempo total destinado à sua resolução, atingiram uma porcentagem que variou de 31,9 \%, na região Norte, até 48,3 \%, na região Centro-Oeste.
Há diferença estatisticamente significativa entre a região Norte e as demais regiões.
Considerando -se o desempenho dos alunos, nota -se ainda que 44,3 \% consideraram a extensão da prova Adequada, no quarto de desempenho inferior, e 45,3 \%, no de melhor desempenho ( quarto superior ).
Nos quartos intermediários, essa proporção foi de 42,8 \%, no segundo quarto, e de 47,7 \%, no terceiro.
Não é discernível uma tendência segundo o desempenho.
No Gráfico 4.10, observa -se que as proporções de estudantes que consideraram a prova Longa ou Muito longa não apresentam uma tendência específica em função dos quartos de desempenho, sendo a maior proporção dessa resposta no segundo quarto de desempenho ( 47,7 \% ), e a menor, nos primeiro e terceiro quartos ( 44,4 \% ).
Pode -se constatar que não há diferença estatisticamente significativa entre a proporção dos quartos de desempenho.
Considerando -se a extensão da prova, o percentual de estudantes que avaliou a prova como Longa ou Muito longa, em relação ao tempo total destinado à sua resolução, foi maior entre os estudantes de CEFET/IF, alcançando 48,9 \%, enquanto entre os estudantes de Faculdades, a Organização Acadêmica com a menor incidência, esse índice atingiu 41,7 \%.
No Gráfico 4.11, é possível observar que não há diferença estatisticamente significativa entre todos os tipos de Organização Acadêmica.
A proporção de presentes à prova que considerou a sua extensão Adequada, em relação ao tempo total destinado à sua resolução, \textbf{esteve} entre 40,0 \%, nos CEFET/IF, e 51,7 \%, nas Faculdades ( ver também Tabela II.6, no Anexo II ).
O percentual de estudantes que considerou a prova Longa ou Muito longa, em relação ao tempo total destinado à sua resolução, foi maior entre os estudantes de cursos de IES Privadas ( 46,6 \% ), não havendo diferença estatisticamente significativa da proporção em relação aos de IES Públicas ( 44,0 \% ).
Já a proporção de presentes à prova que considerou a extensão Adequada foi de 44,1 \% para os de IES Públicas e de 45,6 \% para os de IES Privadas ( ver Gráfico 4.12 e Tabela II.6, no Anexo II ). 4.3 COMPREENSÃO DOS ENUNCIADOS DAS QUESTÕES 4.3.1 COMPONENTE DE FORMAÇÃO GERAL Em relação aos enunciados das questões do Componente de Formação Geral ( Questão 4 ), as opiniões foram positivas, já que 74,6 \% dos alunos avaliados consideraram todos ou a maioria dos enunciados das questões “ claros e objetivos ” ( Gráfico 4.13, Gráfico 4.14 e, no Anexo II, a Tabela II.7 ).
Na análise regional, a percentagem de estudantes que avaliou que todos ou a maioria dos enunciados das questões do Componente de Formação Geral \textbf{estavam} “ claros e objetivos ” variou de 72,0 \%, na região Sudeste, a 79,6 \%, na região Centro-Oeste.
A análise das percepções dos estudantes sobre a clareza e a objetividade dos enunciados permite afirmar que todos ou a maioria dos enunciados de questões relativas ao Componente de Formação Geral, foram considerados “ claros e objetivos ” para a maior parte dos respondentes ( maior ou igual a 72,0 \%, em todas as regiões e maior ou igual a 62,3 \% para todos os quartos de desempenho ).
Segundo o desempenho, observa -se que a proporção dos que emitiram essa opinião cresce, conforme o desempenho aumenta.
Apenas não há diferenças estatisticamente significativas entre os dois quartos superiores.
No quarto superior, a clareza e a objetividade de todos ou da maioria dos enunciados das questões foram percebidas por 82,8 \%.
A proporção dos presentes que considerou que todos ou a maioria dos enunciados das questões do Componente de Formação Geral \textbf{estavam} “ claros e objetivos ” foi maior entre os estudantes de Faculdades, chegando a 83,3 \%, enquanto entre os de Centros Universitários, a Organização Acadêmica com a menor incidência, esse índice atingiu 71,4 \%.
No Gráfico 4.15, é possível observar que a diferença entre as Universidades e os Centros Universitários é estatisticamente significativa ( ver também Tabela II.8, no Anexo II ).
O percentual de estudantes que considerou que todos ou a maioria dos enunciados das questões do Componente de Formação Geral \textbf{estavam} “ claros e objetivos ” foi maior entre os estudantes de cursos de IES Públicas ( 78,3 \% ), com diferença estatisticamente significativa da proporção em IES Privadas ( 72,4 \% ).
Um maior \textbf{detalhamento} pode ser encontrado no Gráfico 4.16 e na Tabela II.8, no Anexo II. 4.3.2 COMPONENTE DE CONHECIMENTO ESPECÍFICO Em relação aos enunciados das questões do Componente de Conhecimento Específico da prova, para 73,6 \% dos estudantes avaliados da Área de Artes Visuais ( Licenciatura ), a clareza e a objetividade ( Questão 5 ) \textbf{estavam} presentes em todas ou na maioria das questões ( Gráfico 4.17 a Gráfico 4.20, e no Anexo II, as Tabelas II.9 e II.10 ).
A maioria dos estudantes de todas as Grandes Regiões brasileiras considerou “ claros e objetivos ” todos ou a maioria dos enunciados das questões do Componente de Conhecimento Específico da prova, percentual sempre maior ou igual a 70,7 \%.
Apenas a diferença entre as regiões Sudeste e Sul é estatisticamente significativa.
A proporção de estudantes que considerou os enunciados das questões “ claros e objetivos ” apresenta uma tendência crescente em relação ao aumento de desempenho : mais elevada no quarto superior ( 81,5 \% ), se comparada ao quarto inferior de desempenho ( 61,7 \% ).
Apenas a diferença entre os dois quartos superiores de desempenho não é estatisticamente significativa.
A proporção dos presentes que considerou que todos ou a maioria dos enunciados das questões do Componente de Conhecimento Específico \textbf{estavam} “ claros e objetivos ” foi maior entre os estudantes de Faculdades, atingindo 80,0 \%.
Entre os estudantes de Centro Universitários, a Organização Acadêmica com a menor incidência, essa proporção atingiu 72,3 \%.
No Gráfico 4.19, é possível observar que nenhuma das diferenças entre os tipos de Organização Acadêmica é estatisticamente significativa ( ver também Tabela II.10, no Anexo II ).
O percentual de estudantes que considerou que todos ou a maioria dos enunciados das questões do Componente de Conhecimento Específico \textbf{estavam} “ claros e objetivos ” foi maior entre os estudantes de cursos de IES Públicas ( 75,8 \% ), sem diferença estatisticamente significativa da proporção em IES Privadas ( 72,3 \% ).
Ver Gráfico 4.20 e Tabela II.10, no Anexo II, para maior \textbf{detalhamento}. 4.4 SUFICIÊNCIA DAS INFORMAÇÕES/INSTRUÇÕES FORNECIDAS Ao se avaliarem as informações/instruções fornecidas para a resolução das questões ( Questão 6 ), constatou -se que 81,8 \% dos respondentes da Área de Artes Visuais ( Licenciatura ) de todo o Brasil afirmaram que estas eram até excessivas ou suficientes em todas ou na maioria das questões ( Gráfico 4.21, Gráfico 4.22 e, no Anexo II, a Tabela II.11 ).
Quanto à distribuição de respondentes pelas Grandes Regiões, observa -se que a proporção de estudantes que considerou as informações/instruções fornecidas até excessivas ou suficientes em todas ou na maioria das questões foi sempre superior ou igual a 79,0 \%, chegando a 84,9 \%, na região Centro-Oeste.
Não existem diferenças estatisticamente significativas entre as regiões.
Levando -se em conta o desempenho dos participantes, observa -se uma tendência crescente das proporções de participantes que avaliaram as informações/instruções como até excessivas ou suficientes em todas ou na maioria das questões, não havendo diferença estatisticamente significativa apenas entre o terceiro e o quarto superior de desempenho, conforme mostrado no Gráfico 4.22.
O percentual foi mais elevado no quarto superior ( 90,7 \% ), percentual superior à média nacional ( 81,8 \% ).
Já no quarto inferior, a suficiência das informações/instruções declarada como até excessiva, em todas ou na maioria das questões foi percebida por 70,1 \% dos respondentes.
Considerando -se as informações/instruções fornecidas para a resolução das questões ( Questão 6 ), a proporção que considerou que estas eram até excessivas ou suficientes em todas ou na maioria das questões foi maior entre os estudantes de Universidades, totalizando 83,9 \%.
Os Centros Universitários foram a Organização Acadêmica em que houve a menor proporção, 77,9 \%.
No Gráfico 4.23, é possível observar que apenas a diferença entre as Universidades e os Centros Universitários é estatisticamente significativa ( ver também Tabela II.12, no Anexo II ).
O percentual de estudantes que considerou que as informações/instruções fornecidas para a resolução das questões eram até excessivas ou suficientes em todas ou na maioria das questões foi maior para os estudantes de cursos de IES Públicas ( 86,4 \% ), havendo diferença estatisticamente significativa da proporção em IES Privadas ( 79,0 \% ), ver Gráfico 4.24 e Tabela II.12, no Anexo II. 4.5 DIFICULDADE ENCONTRADA AO RESPONDER À PROVA Indagados sobre as dificuldades com as quais se depararam ao responder à prova ( Questão 7 ), 15,0 \% dos estudantes apontaram o Desconhecimento do conteúdo.
Para 44,4 \%, a Forma diferente de abordagem do conteúdo foi indicada como dificuldade.
Já a Falta de motivação para fazer a prova foi a dificuldade apontada por 18,2 \% dos respondentes.
Considerando -se todo o Brasil, 14,9 \% dos respondentes afirmaram que não tiveram qualquer tipo de dificuldade para responder à prova ( Tabelas II.13 e II.14, no Anexo II ).
Nos Gráficos 4.25 a 4.28, são apresentados os percentuais de estudantes que apontaram o Desconhecimento do conteúdo como “ dificuldade ao responder à prova ”.
Na análise por Grandes Regiões, o percentual de inscritos e presentes que apontaram o Desconhecimento do conteúdo como “ dificuldade ao responder à prova ” não superou 18,7 \%.
Os percentuais variaram de 8,6 \%, na região Centro-Oeste, a 18,7 \%, na região Norte.
A Forma diferente de abordagem do conteúdo foi a escolha modal dos estudantes, com percentuais que variaram de 39,6 \% ( região Nordeste ) a 49,7 \% ( região Centro-Oeste ).
O percentual de alunos que citou a Falta de motivação como dificuldade variou de 14,4 \% ( região Norte ) a 24,0 \% ( região Nordeste ).
Os que declararam não ter qualquer tipo de dificuldade para responder à prova variaram de 12,6 \%, na região Centro-Oeste, a 15,8 \%, na região Norte.
Em relação aos quartos de desempenho, o Desconhecimento do conteúdo foi a opção escolhida por 17,8 \% dos estudantes do quarto inferior e por 10,2 \% do quarto superior, caracterizando uma tendência decrescente a partir do segundo quarto.
São observadas diferenças estatisticamente significativas entre os dois quartos inferiores e os dois quartos superiores de desempenho.
A alternativa modal para os alunos, quando agregados pelos quartos de desempenho, para a dificuldade encontrada, foi causada pela Forma diferente de abordagem do conteúdo : 46,0 \%, no primeiro quarto, e 41,7 \% no quarto superior, escolheram essa alternativa.
Na análise por Organização Acadêmica, o percentual de inscritos e presentes que apontou o Desconhecimento do conteúdo como “ dificuldade ao responder à prova ” não superou 20,2 \%.
Os percentuais variaram de 12,1 \%, nas Universidades, a 20,2 \%, nos Centros Universitários.
A Forma diferente de abordagem do conteúdo foi a escolha modal dos estudantes, com exceção daqueles de Faculdades, com moda na Falta de motivação para fazer a prova ( 35,0 \% ).
Para os demais tipos, os percentuais variaram de 42,4 \% ( Universidades ) a 51,1 \% ( CEFET/IF ).
O percentual de alunos que citou a Falta de motivação como dificuldade variou de 13,5 \% ( Centros Universitários ) a 35,0 \% ( Faculdades ).
Os que declararam não ter qualquer tipo de dificuldade para responder à prova variaram de 8,9 \%, nos CEFET/IF, a 20,0 \%, nas Faculdades.
No Gráfico 4.27, é possível observar que há diferença estatisticamente significativa entre os alunos de Universidades e de Centros Universitários com respeito à alternativa Desconhecimento do conteúdo ( ver também Tabela II.14, no Anexo II ).
Em relação às Categorias Administrativas, o Desconhecimento do conteúdo foi a opção escolhida por 9,8 \% dos estudantes nas IES Públicas e por 18,1 \% nas IES Privadas.
É observada uma diferença estatisticamente significativa entre essas duas categorias.
A alternativa modal para os alunos, quando agregados pelas Categorias Administrativas, para a dificuldade encontrada, foi causada pela Forma diferente de abordagem do conteúdo : 40,1 \%, para IES Públicas e 47,0 \%, para as IES Privadas ( ver Gráfico 4.28 e Tabela II.14, no Anexo II ). 4.6 CONTEÚDOS DAS QUESTÕES OBJETIVAS DA PROVA Ao se analisarem os conteúdos das questões objetivas da prova ( Questão 8 ), verifica -se que 13,0 \% dos estudantes avaliados afirmaram que não estudou ainda a maioria desses conteúdos ( Gráficos 4.29 a Gráfico 4.32, e as Tabelas II.15 e II.16, no Anexo II ).
A maioria ( 60,0 \% ) afirmou ter estudado e aprendido muitos ou todos os conteúdos.
Na análise por Grande Região, a proporção de respondentes que escolheu a opção não estudou ainda a maioria desses conteúdos foi maior do que a média nacional nas regiões Norte ( 17,1 \% ), Sudeste ( 14,3 \% ) e Sul ( 13,4 \% ).
Nas regiões Nordeste e Centro-Oeste, a situação foi inversa, com proporções mais baixas, correspondendo a, respectivamente, 6,6 \% e 7,3 \%.
São observadas diferenças estatisticamente significativas entre a região Nordeste e as regiões Norte, Sudeste e Sul, e também entre a região Centro-Oeste e as regiões Sudeste e Sul.
Em todas as regiões, cerca de 50 \% ou mais dos presentes afirmou ter estudado e aprendido muitos ou todos os conteúdos, com proporções variando entre 47,1 \%, na região Norte, a 76,2 \%, na região Nordeste.
Considerando -se separadamente as opiniões de estudantes dos quatro quartos de desempenho, observa -se que, no quarto inferior, 23,4 \% ofereceram como resposta que não estudou ainda a maioria desses conteúdos, sendo 3,9 \% os do quarto superior com a mesma resposta.
As proporções são decrescentes com o desempenho.
Apenas a diferença entre os alunos que optaram por esse motivo de dificuldade nos dois quartos inferiores de desempenho não é estatisticamente significativa.
Tendo -se em conta o quarto superior, 85,9 \% dos alunos afirmaram ter estudado e aprendido muitos ou todos os conteúdos.
No outro extremo, no primeiro quarto, 35,4 \% optaram pelas mesmas categorias.
As proporções são crescentes com o desempenho.
Na análise por Organização Acadêmica, a proporção de respondentes que escolheu a opção não estudou ainda a maioria desses conteúdos não superou 18,5 \%.
Observa -se que as proporções nas Universidades e CEFET/IF, com 10,1 \% e 7,0 \% respectivamente, foram menores do que a média nacional ( 13,0 \% ).
Nos Centros Universitários e nas Faculdades, a situação foi inversa : a proporção foi mais alta, com 18,5 \% e 16,7 \% respectivamente.
No Gráfico 4.31, é possível observar que as diferenças entre os Centros Universitários e o CEFET/IF e Universidades são estatisticamente significativas.
Em todos os tipos de Organização Acadêmica, cerca de 50 \% ou mais dos presentes afirmou ter estudado e aprendido muitos ou todos os conteúdos, com proporções variando entre 49,4 \%, nos Centros Universitários, a 72,1 \%, nos CEFET/IF ( ver também Tabela II.16, no Anexo II ).
Em relação às Categorias Administrativas, observa -se que, nas IES Públicas, 6,4 \% ofereceram como resposta que não estudou ainda a maioria desses conteúdos, sendo 17,1 \% nas IES Privadas com a mesma resposta.
A diferença entre os alunos que optaram por esse motivo de dificuldade nas duas categorias é estatisticamente significativa.
Considerando -se as IES Públicas, 74,2 \% dos alunos afirmaram ter estudado e aprendido muitos ou todos os conteúdos, sendo de 51,3 \% a proporção dos que escolheram essa resposta nas IES Privadas, com uma diferença estatisticamente significativa ( ver Gráfico 4.32 e Tabela II.16, no Anexo II ). 4.7 TEMPO GASTO PARA CONCLUIR A PROVA Ao responderem sobre o tempo de conclusão da prova ( Questão 9 ), mais de três quartos dos estudantes ( 75,9 \% ) afirmou ter gasto Entre duas e quatro horas ( Gráfico 4.33 a Gráfico 4.36 e, no Anexo II, as Tabelas II.17 e II.18 ).
Considerando -se as cinco Grandes Regiões brasileiras, os percentuais dos que utilizaram Entre duas e quatro horas para finalizar a prova nas regiões Norte ( 72,1 \% ), Sudeste ( 75,7 \% ) e Centro-Oeste ( 75,5 \% ) foram inferiores ao percentual nacional.
Nas demais Grandes Regiões, os percentuais de alunos que dispensaram Entre duas e quatro horas para concluir a prova ficaram acima da média nacional : 77,4 \% na região Nordeste e 76,2 \%, na região Sul, como mostrado no Gráfico 4.33.
Não há diferenças estatisticamente significativas entre as regiões.
Uma vez analisadas as alternativas escolhidas pelos estudantes que se situam nos diferentes quartos de desempenho, observa -se uma tendência crescente com o desempenho : uma proporção maior de participantes no quarto superior declarou ter gasto Entre duas e quatro horas para concluir a prova quando comparada com a proporção de participantes nos quartos inferiores.
Há diferenças estatisticamente significativas apenas entre o quarto inferior e os demais quartos de desempenho.
Na análise por Organização Acadêmica, os percentuais dos que utilizaram Entre duas e quatro horas para finalizar a prova nos Centros Universitários ( 77,6 \% ), nas Faculdades ( 83,3 \% ) e nos CEFET/IF ( 82,2 \% ) foram superiores ao percentual nacional ( 75,9 \% ).
Nas Universidades, tal percentual foi de 74,6 \%, e, como mostrado no Gráfico 4.35, não há diferença estatisticamente significativa entre as proporções dos tipos de Organização Acadêmica.
Em relação às Categorias Administrativas, observa -se que, nas IES Públicas, 75,4 \% declararam ter gasto Entre duas e quatro horas para concluir a prova, sendo 76,1 \% nas IES Privadas com a mesma resposta.
A diferença entre os alunos que optaram por esse motivo de dificuldade nas duas categorias não é estatisticamente significativa ( ver Gráfico 4.36 e Tabela II.18, no Anexo II ).

% matched lemmas: cada, certo, detalhamento, estar, lembrar
\end{textsample}
